% GENERATED FILE. Edit canonical lessons, then run npm run generate:books.
% canonical-source-hash: fnv1a64:a27cd68fa795a93b
% canonical-lessons: TA-C455-cancalam, TA-C455-ulaiccal, TA-C455-atankam, TA-C455-ankalay, TA-C455-tavi

\chapter{Worry and Unease}
\label{ch:ta-worry-and-unease}
\hlchaptermodality{455}

\begin{chapteropening}
\textbf{By the end of this chapter:} \emph{I can talk about unease, anguish and fretting.}

Complete the last lesson of chapter 455: I can talk about unease, anguish and fretting.
\end{chapteropening}

\section[\texorpdfstring{\ta{சஞ்சலம்} (cañcalam)}{சஞ்சலம் (cañcalam)}]{\ta{சஞ்சலம்} (cañcalam) — unease; wavering}

\label{lesson:TA-C455-cancalam}



\begin{quote}
Before the new one: say the Tamil for to be startled, then the Tamil for to tremble.
\end{quote}

\subsection*{You'll want to know: \ta{சஞ்சலம்}}
\textbf{\ta{சஞ்சலம்}} — \emph{cañcalam} — \textquotedblleft{}unease; wavering\textquotedblright{}.

\begin{quote}
\ta{முடிவு} \ta{செய்ய} \ta{முடியாமல்} \ta{அவர்} \ta{சஞ்சலத்தில்} \ta{இருந்தார்}.
\end{quote}

\begin{quote}
\emph{muṭivu ceyya muṭiyāmal avar cañcalattil iruntār.}
\end{quote}

\begin{quote}
Unable to decide, he was uneasy.
\end{quote}

The oblique -\ta{அத்து}-: \ta{சஞ்சலத்தில்}, in a state of unease.

\begin{cousinweb}
\textbf{\ta{சஞ்சலம்}} (\emph{cañcalam}) is borrowed from Sanskrit \emph{cañcala}, \textquotedblleft{}quivering, unsteady\textquotedblright{} (CDIAL 4568), with the Tamil noun ending -\ta{அம்}.
\end{cousinweb}

\begin{grammarlens}[title={What you've built}]
One more word for describing what you lived and felt.
\end{grammarlens}

\subsection*{Your turn}
\begin{itemize}
\raggedright
  \item \emph{Say it:} \emph{cañcalam}
  \item \emph{Say it:} \emph{cañcalam}, once more
  \item \emph{Say it:} \emph{muṭivu ceyya muṭiyāmal avar cañcalattil iruntār.} — Unable to decide, he was uneasy.
  \item \emph{Recall:} say the Tamil for to be startled, then the Tamil for to tremble, then say \emph{cañcalam} again
\end{itemize}

\begin{tcolorbox}[breakable,colback=teal!4,colframe=teal!35!black,title={Before you move on}]
What does \ta{சஞ்சலம்} mean? (\textquotedblleft{}Unease; wavering\textquotedblright{}.) Say it once more.
\end{tcolorbox}
\section[\texorpdfstring{\ta{உளைச்சல்} (uḷaiccal)}{உளைச்சல் (uḷaiccal)}]{\ta{உளைச்சல்} (uḷaiccal) — anguish; mental strain}

\label{lesson:TA-C455-ulaiccal}



\begin{quote}
Before the new one: say the Tamil for to tremble, then the Tamil for unease.
\end{quote}

\subsection*{You'll want to know: \ta{உளைச்சல்}}
\textbf{\ta{உளைச்சல்}} — \emph{uḷaiccal} — \textquotedblleft{}anguish; mental strain\textquotedblright{}.

\begin{quote}
\ta{வேலையை} \ta{இழந்த} \ta{பிறகு} \ta{அவருக்கு} \ta{மன} \ta{உளைச்சல்} \ta{ஏற்பட்டது}.
\end{quote}

\begin{quote}
\emph{vēlaiyai iḻanta piṟaku avarukku maṉa uḷaiccal ēṟpaṭṭatu.}
\end{quote}

\begin{quote}
After losing his job he suffered mental anguish.
\end{quote}

Mostly in \ta{மன} \ta{உளைச்சல்}, mental anguish, with \ta{மனம்} \textquotedblleft{}mind\textquotedblright{}.

\begin{cousinweb}
\textbf{\ta{உளைச்சல்}} (\emph{uḷaiccal}) is native Dravidian, built on \ta{உளை} \emph{uḷai}, \textquotedblleft{}to ache, to suffer pain\textquotedblright{}, which DEDR 2702 lists with a query.
\end{cousinweb}

\begin{grammarlens}[title={What you've built}]
One more word for describing what you lived and felt.
\end{grammarlens}

\subsection*{Your turn}
\begin{itemize}
\raggedright
  \item \emph{Say it:} \emph{uḷaiccal}
  \item \emph{Say it:} \emph{uḷaiccal}, once more
  \item \emph{Say it:} \emph{vēlaiyai iḻanta piṟaku avarukku maṉa uḷaiccal ēṟpaṭṭatu.} — After losing his job he suffered mental anguish.
  \item \emph{Recall:} say the Tamil for to tremble, then the Tamil for unease, then say \emph{uḷaiccal} again
\end{itemize}

\begin{tcolorbox}[breakable,colback=teal!4,colframe=teal!35!black,title={Before you move on}]
What does \ta{உளைச்சல்} mean? (\textquotedblleft{}Anguish; mental strain\textquotedblright{}.) Say it once more.
\end{tcolorbox}
\section[\texorpdfstring{\ta{ஆதங்கம்} (ātaṅkam)}{ஆதங்கம் (ātaṅkam)}]{\ta{ஆதங்கம்} (ātaṅkam) — frustration; anguished regret}

\label{lesson:TA-C455-atankam}



\begin{quote}
Before the new one: say the Tamil for unease, then the Tamil for anguish.
\end{quote}

\subsection*{You'll want to know: \ta{ஆதங்கம்}}
\textbf{\ta{ஆதங்கம்}} — \emph{ātaṅkam} — \textquotedblleft{}frustration; anguished regret\textquotedblright{}.

\begin{quote}
\ta{மகன்} \ta{படிக்கவில்லை} \ta{என்று} \ta{அப்பாவுக்கு} \ta{ஆதங்கம்}.
\end{quote}

\begin{quote}
\emph{makaṉ paṭikkavillai eṉṟu appāvukku ātaṅkam.}
\end{quote}

\begin{quote}
The father is anguished that his son did not study.
\end{quote}

A frustrated wish for someone else: … \ta{என்று} \ta{ஆதங்கம்}, anguish that ….

\begin{cousinweb}
\textbf{\ta{ஆதங்கம்}} (\emph{ātaṅkam}) is borrowed from Sanskrit \emph{ātaṅka}, \textquotedblleft{}disease; affliction, anxiety\textquotedblright{} (Monier-Williams).
\end{cousinweb}

\begin{grammarlens}[title={What you've built}]
One more word for describing what you lived and felt.
\end{grammarlens}

\subsection*{Your turn}
\begin{itemize}
\raggedright
  \item \emph{Say it:} \emph{ātaṅkam}
  \item \emph{Say it:} \emph{ātaṅkam}, once more
  \item \emph{Say it:} \emph{makaṉ paṭikkavillai eṉṟu appāvukku ātaṅkam.} — The father is anguished that his son did not study.
  \item \emph{Recall:} say the Tamil for unease, then the Tamil for anguish, then say \emph{ātaṅkam} again
\end{itemize}

\begin{tcolorbox}[breakable,colback=teal!4,colframe=teal!35!black,title={Before you move on}]
What does \ta{ஆதங்கம்} mean? (\textquotedblleft{}Frustration; anguished regret\textquotedblright{}.) Say it once more.
\end{tcolorbox}
\section[\texorpdfstring{\ta{அங்கலாய்} (aṅkalāy)}{அங்கலாய் (aṅkalāy)}]{\ta{அங்கலாய்} (aṅkalāy) — to bemoan; to fret}

\label{lesson:TA-C455-ankalay}



\begin{quote}
Before the new one: say the Tamil for anguish, then the Tamil for frustration.
\end{quote}

\subsection*{You'll want to know: \ta{அங்கலாய்}}
\textbf{\ta{அங்கலாய்}} — \emph{aṅkalāy} — \textquotedblleft{}to bemoan; to fret\textquotedblright{}.

\begin{quote}
\ta{விலை} \ta{ஏறியதைப்} \ta{பற்றிப்} \ta{பாட்டி} \ta{அங்கலாய்த்தாள்}.
\end{quote}

\begin{quote}
\emph{vilai ēṟiyataip paṟṟip pāṭṭi aṅkalāyttāḷ.}
\end{quote}

\begin{quote}
Grandmother fretted about prices going up.
\end{quote}

Past (I …): \ta{அங்கலாய்த்தேன்}, \emph{aṅkalāyttēṉ}.

What one frets about takes \ta{பற்றி} \textquotedblleft{}about\textquotedblright{}: \ta{அதைப்} \ta{பற்றி} \ta{அங்கலாய்த்தாள்}.

\begin{cousinweb}
\textbf{\ta{அங்கலாய்}} (\emph{aṅkalāy}) is native Dravidian (DEDR 31); the entry also gives \ta{அங்கலாய்ப்பு} \emph{aṅkalāyppu}, \textquotedblleft{}disquiet\textquotedblright{}.
\end{cousinweb}

\begin{grammarlens}[title={What you've built}]
One more word for describing what you lived and felt.
\end{grammarlens}

\subsection*{Your turn}
\begin{itemize}
\raggedright
  \item \emph{Say it:} \emph{aṅkalāy}
  \item \emph{Say it:} \emph{aṅkalāy}, once more
  \item \emph{Say it:} \emph{vilai ēṟiyataip paṟṟip pāṭṭi aṅkalāyttāḷ.} — Grandmother fretted about prices going up.
  \item \emph{Recall:} say the Tamil for anguish, then the Tamil for frustration, then say \emph{aṅkalāy} again
\end{itemize}

\begin{tcolorbox}[breakable,colback=teal!4,colframe=teal!35!black,title={Before you move on}]
What does \ta{அங்கலாய்} mean? (\textquotedblleft{}To bemoan; to fret\textquotedblright{}.) Say it once more.
\end{tcolorbox}
\section[\texorpdfstring{\ta{தவி} (tavi)}{தவி (tavi)}]{\ta{தவி} (tavi) — to suffer; to struggle}

\label{lesson:TA-C455-tavi}



\begin{quote}
Before the new one: say the Tamil for frustration, then the Tamil for to bemoan.
\end{quote}

\subsection*{You'll want to know: \ta{தவி}}
\textbf{\ta{தவி}} — \emph{tavi} — \textquotedblleft{}to suffer; to struggle\textquotedblright{}.

\begin{quote}
\ta{தண்ணீர்} \ta{இல்லாமல்} \ta{மக்கள்} \ta{தவித்தார்கள்}.
\end{quote}

\begin{quote}
\emph{taṇṇīr illāmal makkaḷ tavittārkaḷ.}
\end{quote}

\begin{quote}
People suffered without water.
\end{quote}

Past (I …): \ta{தவித்தேன்}, \emph{tavittēṉ}.

Often with -\ta{ஆமல்} \textquotedblleft{}without\textquotedblright{}: \ta{பணம்} \ta{இல்லாமல்} \ta{தவித்தேன்}, I struggled without money.

\begin{cousinweb}
\textbf{\ta{தவி}} (\emph{tavi}) is a native Tamil verb, but the origin is not settled: DEDR gives it no entry of its own that we can point to.
\end{cousinweb}

\begin{grammarlens}[title={What you've built}]
One more word for describing what you lived and felt.
\end{grammarlens}

\subsection*{Your turn}
\begin{itemize}
\raggedright
  \item \emph{Say it:} \emph{tavi}
  \item \emph{Say it:} \emph{tavi}, once more
  \item \emph{Say it:} \emph{taṇṇīr illāmal makkaḷ tavittārkaḷ.} — People suffered without water.
  \item \emph{Recall:} say the Tamil for frustration, then the Tamil for to bemoan, then say \emph{tavi} again
\end{itemize}

\begin{tcolorbox}[breakable,colback=teal!4,colframe=teal!35!black,title={Before you move on}]
What does \ta{தவி} mean? (\textquotedblleft{}To suffer; to struggle\textquotedblright{}.) Say it once more.
\end{tcolorbox}
